%!TEX root = ../Polito PhD thesis template.tex
% ******************************* Chapter 9 ************************************

\chapter{Conclusions and Perspectives}\label{ch:conclusions}

This thesis has dealt with the data that a neutronic analysis consults in
place of the Monte Carlo reference it cannot afford to repeat: the multi-group
constants of a deterministic solver for
two lead-cooled fast reactors, and a database of fission matrices for a
heat-pipe micro reactor. The conclusions of each Part are gathered here, and
the chapter closes with the developments that the work leaves open.

\section{Multi-group data for lead-cooled fast reactors}

For the Westinghouse LFR, the group constants were generated with Serpent~2
either from one model of the whole core or from mini-core models built region
by region, and fed to a diffusion solver based on finite elements. The mini-core route proves the better investment, in accuracy
and in computing time, whenever several configurations of the core have to be
analysed, since a full-core model has to be rerun for each of them, at a cost
that grows when low statistical errors are wanted in the non-multiplying
regions. The superhomogenization correction of the regions responsible for the
largest discrepancies brings the multiplication factor and the power density
distribution of the deterministic model close to the reference, at a cost of
seconds per supercell, and it does so more effectively where the system is more
heterogeneous. The kinetic parameters respond to the correction without a
direct benefit, which follows from what the correction preserves: reaction
rates, the product of a cross section and the flux, and not bilinear quantities
weighted with the adjoint. Fast reactors proved favourable ground for the
method, since the relatively uniform flux, the reduced anisotropy and the
smooth energy dependence of the cross sections limit the condensation errors
and let scalar factors preserve the global quantities.

For the ALFRED core the object of the optimization was the grid itself, the
set of energy boundaries on which the constants are condensed, and the dependence of the result on
the random seed, observed in the work this chapter built on, was measured on
one thousand independent searches instead of three. The number was made
affordable by a neural metamodel trained on $91\,723$ full-core evaluations of
FRENETIC, which replaced the solver inside the search alone; every structure
the searches returned was recomputed with the solver, so that the conclusions
rest on it and not on the metamodel. The dependence is severe: a single search
returns a structure typically forty-five per cent above the best on the joint
objective, three searches twenty per cent, and a hundred two per cent, so that
the gain from repeating the search lies precisely in the range that the
reference solver cannot reach. The thousand searches converged to 747 distinct
structures with no common skeleton; the regularities they display are an
exclusive choice between two boundaries of the fast region and a placement of
the boundaries that follows the neutron importance rather than the flux, a
criterion never imposed and found a posteriori. The best structure has fifteen groups and improves on
the three obtained with the solver inside the loop, on the structures built by
uniform spacing rules and on the finest structure the collapsing procedure
admits, the last by five orders of magnitude on the joint objective. That last
comparison is the one that explains the result: a structure that condenses
almost nothing leaves the error of the deterministic model intact, and the
optimized structure does better because its condensation error cancels that
model error. The quality of the best structures is therefore bound to the core,
to the solver and to the objective on which it was obtained, and what carries
over to another case is the procedure, whose cost was repaid after nine
searches and which delivered analyses worth of the order of $390\,000$
core-hours of the reference solver for twenty-nine.

\section{Fission matrices for heat-pipe micro reactors}

The heat-pipe micro reactor was studied on the model of the Virtual Test Bed.
Its criticality search rests on a database of 136 fission matrices, tallied at
eight fuel temperatures and seventeen angles of the control drums, and on a
prescribed relation between power and fuel temperature that takes the place of
a thermal solver; the density of the database was not assumed but established
by a convergence study. With fresh fuel at nominal power one search takes about
a second on a single core, its uncertainty quantification included, and a
reference Monte Carlo calculation at the state it returns gives a
multiplication factor $-56\pm16$~pcm away from the value sought, a deviation
that is the residual of the interpolation of the database. Burnup was brought
in without enlarging the database: a correction ratio, tallied once for each
snapshot of the cycle at a single reference state, multiplies the interpolated
matrix, and the cycle is followed quasi-statically, with drum angle,
temperatures, fission source and cell-wise burnup advancing together. Over
sixty-six months the states so obtained have a multiplication factor between
$-59$ and $-5$~pcm from unity when recomputed by Monte Carlo, with the same
bias found at beginning of life and no drift with burnup, and an assembly-wise
fission source reproduced to about one per cent. The three assumptions on
which the factorization rests were each tested on a dedicated campaign and
found to cost no more than that residual, and seven snapshots along the cycle
proved sufficient.

The isothermal database cannot represent a reflector colder than the core, nor
a temperature that changes sharply from one cell to the next, and each of these
effects received a correction ratio with its own reference case and its own
weight. For the reflector, the correction takes the deviation of the
multiplication factor on a uniformly hot core from several hundred pcm down to
a few tens, and to less than the statistical uncertainty of the reference once
the reference matrix is tallied at the angle of the drums, at the price of one
further calculation per angle; on fields that vary across the core it removes
almost all of the discrepancy when the variation is gradual and only a part of
it when the temperature jumps from ring to ring. For the gradients inside the
core, the correction reinstates the influence that the neighbouring cells exert
on a coefficient, and over twelve temperature fields it lowers the mean
absolute deviation of the multiplication factor by a quarter and that of the
fission source by as much, the gain being largest on fields that are both
sharp and extended over several cells. The statistical uncertainty of
everything obtained from an interpolated matrix was also examined. The
coefficients that share a source column turn out to be negatively correlated,
because the histories of that source are divided among destinations that
exclude one another, and the correlation follows in closed form from the mean
coefficients and the multiplicity alone; checked against 150 independent
databases, the closed form gives the uncertainty of the multiplication factor
to within five per cent, whereas the uncorrelated assumption used elsewhere in
the Part overestimates it, so that the bands reported for the criticality
search err on the safe side.

\section{Perspectives}

Each of the works leaves a direction open, and they are stated as the chapters
stated them. On the multi-group side, the mini-core route to the generation of
the constants lends itself to sensitivity and uncertainty studies on the nuclear data in
transient calculations, and the reproduction of the kinetic
parameters under superhomogenization correction, which preserves reaction rates
and not adjoint-weighted quantities, remains to be improved. The neural
metamodel is still limited by the amount of reference data, and the evaluations
that would help it most are those where the target is steepest, in the
neighbourhood of the optimized structures, which a campaign of the reference
solver aimed at that region could provide; the network could also receive, in
addition to the boundaries of a structure, quantities that carry physical
content, such as the lethargy width of each condensed group or the variation of
the reference spectrum inside it, a possibility that a first attempt left open.

On the fission matrix side, what limits the criticality search is the
interpolation of the original database rather than the burnup correction, and
the prescribed relation between power and fuel temperature can give way to a
thermal response computed by a conduction solver, for which the workflow
already provides the interface. The two temperature corrections were built on
different models and used one at a time; using them together requires
calibrating both on one and the same model, together with a weight that avoids
counting the same effect twice in the outermost ring of assemblies, where a hot
core and a sharp gradient overlap. Where the temperature falls from ring to ring, a reference
matrix tallied zone by zone would address the part of the discrepancy that the
reflector correction leaves behind.
